\emph{Foundation model} (FM) keeps its original operational sense~\cite{bommasani2021opportunities,moor2023generalist}:
pre-training at scale on broad medical data, and reuse, frozen or adapted, across several downstream
tasks. Three decisions are kept apart in the eligibility coding. \emph{Foundation-model status} is the
definition above, operationalised for sampling as pre-training on at least $10^{5}$ images or image--text
pairs, $10^{3}$ whole slides or $10^{4}$ volumetric scans from more than one anatomical site or finding
class, and reuse across at least three reported tasks. The three thresholds sit below the smallest
pre-training corpus of each kind in Table~\ref{tab:models}, so that they exclude task-specific models
rather than small foundation models. They compare information units that are not equivalent, since an
image, a slide and a volume differ by orders of magnitude in content, so they are sampling boundaries
rather than part of the definition, and the conclusions do not depend on them: excluding the six
borderline studies changes none of the statements of Table~\ref{tab:statements}; excluding also the
seventeen studies admitted by the volumetric threshold empties the brain-MRI row of the site finding
(F3c) and narrows the modality span of two others (F2 and F8a$^{\prime\prime}$) without reversing
either. Medical adaptations of general foundation models qualify, as do their component encoders.
\emph{Representation access} is whether the study could reach embeddings or activations at all. It is a
requirement of this review rather than a property of the model, and it biases the corpus towards openly
released weights, since API-only models enter only through embedding-level analyses and are marked
borderline. \emph{Claim eligibility} is whether the study makes a representation-level claim about such
a model (\S\ref{sec:method}). One family carries a stricter version of it: a geometry study enters only
if it connects its statistic to decodability, retrieval, transfer or shortcut behaviour
(\S\ref{sec:geometry}). The structural family is therefore reviewed as \emph{task-linked geometry}, and
purely descriptive or negative geometry results are under-represented here relative to the literature. A
study enters if at least one analysed model passes the first two decisions and the study passes the
third; ineligible comparators are named in the evidence tables and support no finding alone. The
reviewed models fall into three architectural templates (Fig.~\ref{fig:overview}, Table~\ref{tab:models}).

\textbf{Unimodal vision encoders} are self-supervised on images alone and reused frozen. They include chest-radiograph
encoders~\cite{sellergren2022simplified,perezgarcia2024scalable}, pathology tile encoders pre-trained on
$10^{8}$--$10^{9}$ whole-slide image (WSI)
tiles~\cite{chen2024generalpurpose,vorontsov2024model,xu2024wholeslide,zimmermann2024virchow2,filiot2024phikon},
retinal~\cite{zhou2023model} and volumetric CT~\cite{blankemeier2024merlin} encoders, and slide-level models
that add weakly supervised stages with interpretable attention over frozen
tiles~\cite{wang2024pathology,ding2024whole}. Promptable segmenters
qualify~\cite{ma2023anything,zhao2024foundation} but are analysed only as encoders; their mask decoders
are unstudied (\S\ref{sec:discovery}). Frozen and public, these encoders carry most of the decoding and
discovery studies (\S\ref{sec:methods}).

\textbf{Dual encoders} align image and text encoders on image--report or caption pairs for zero-shot
classification and retrieval. Chest radiography has the deepest line, from ConVIRT and GLoRIA to CheXzero,
MedCLIP, BioViL and KAD~\cite{zhang2020visual,huang2021gloria,tiu2022detection,wang2022medclip,boecking2022most,zhang2023visuallanguage},
and the recipe transfers wherever paired text exists: BiomedCLIP, CONCH and PLIP, and Mammo-CLIP,
EchoCLIP and
CT-CLIP~\cite{zhang2023biomedclip,lu2024visuallanguage,huang2023visual,ghosh2024mammoclip,christensen2024vision,hamamci2024ctrate}.
Their joint spaces carry the modality-gap and concept-scoring studies of \S\ref{sec:geometry} and
\S\ref{sec:decoding}.

\textbf{Multimodal large language models} feed encoder features through a projector, resampler or
cross-attention connector into a text-generating LLM. The reviewed models range from adapted general
assistants to bespoke medical stacks; Table~\ref{tab:models} gives their connectors, data and weight
availability~\cite{li2023llavamed,moor2023medflamingo,bannur2024maira2,chen2024model,chaves2024clinically,wu2023model,chen2024huatuogptvision,tu2023biomedical,sellergren2025medgemma,team2025lingshu}.
The inherited vision encoder is usually frozen, so its contents remain available even if they are not
used. The connector does more than change the medium: it transforms the encoder representation, and may
compress or discard part of it, before the language model reads it (\S\ref{sec:loci}).

\begin{table*}[!tb]
\centering\footnotesize\renewcommand{\arraystretch}{0.95}
\caption{Representative medical foundation models in the reviewed corpus, by architectural template (CXR: chest radiograph): the models
that recur across the coded records, plus the first model of each template, one to three per modality.
Models analysed by a single core study are named where that study is discussed.
\emph{Connector}: MLP (linear or multilayer projector), resampler (query-based Perceiver or Q-Former),
x-attn (gated cross-attention). \emph{Year}: first public appearance, which may precede the cited version.
\emph{Weights}: released (\cmark), gated (\pmark) or not released (\xmark) at the time of writing.
A dash marks a value the source does not state.}
\label{tab:models}
\setlength\tabcolsep{4pt}\begin{tabular}{@{}L{2.8cm}cL{1.9cm}L{2.5cm}L{2.1cm}L{3.4cm}c@{}}
\toprule
Model & Year & Modality & Vision encoder & Connector & Pre-training data & Weights \\
\midrule
\multicolumn{7}{@{}l}{\emph{Unimodal vision encoders (including slide-level and segmentation models)}} \\
CXR Foundation~\cite{sellergren2022simplified} & 2022 & CXR & EfficientNet-L2 (supervised + contrastive) & --- & large multi-site CXR corpora & \xmark\ (API) \\
RAD-DINO~\cite{perezgarcia2024scalable} & 2024 & CXR & ViT (DINOv2 recipe) & --- & image-only CXR & \cmark \\
UNI~\cite{chen2024generalpurpose} & 2024 & pathology & ViT-L (DINOv2) & --- & $>$100M tiles, $>$100k WSIs & \pmark \\
Virchow~\cite{vorontsov2024model} & 2024 & pathology & ViT-H & --- & 1.5M WSIs & \pmark \\
Prov-GigaPath~\cite{xu2024wholeslide} & 2024 & pathology & ViT-g + slide encoder & --- & 1.3B tiles, 171k WSIs & \pmark \\
RETFound~\cite{zhou2023model} & 2023 & retina & ViT-L (MAE) & --- & 1.6M retinal images & \cmark \\
Merlin~\cite{blankemeier2024merlin} & 2024 & CT & 3-D CNN & --- & 15k CTs + EHR codes/reports & \cmark \\
Virchow2~\cite{zimmermann2024virchow2} & 2024 & pathology & ViT-H (DINOv2, mixed magnification) & --- & 3.1M WSIs & \pmark \\
Phikon-v2~\cite{filiot2024phikon} & 2024 & pathology & ViT-L (DINOv2) & --- & 460M public tiles & \cmark \\
CHIEF~\cite{wang2024pathology} & 2024 & pathology (slide) & tile encoder + weakly supervised slide model & --- & 60.5k WSIs & \cmark \\
TITAN~\cite{ding2024whole} & 2024 & pathology (slide) & slide ViT over CONCH tile features & --- & 336k WSIs + reports/captions & \pmark \\
MedSAM~\cite{ma2023anything} & 2024 & multi (segmentation) & ViT-B (SAM) + mask decoder & --- & 1.57M image--mask pairs & \cmark \\
BiomedParse~\cite{zhao2024foundation} & 2024 & multi (segmentation) & image encoder + text-prompted decoder & --- & $>$6M image--mask--text triples & \cmark \\
\midrule
\multicolumn{7}{@{}l}{\emph{Dual encoders (image--text contrastive)}} \\
ConVIRT~\cite{zhang2020visual} & 2020 & CXR & ResNet-50 & --- & image--report pairs & \cmark \\
GLoRIA~\cite{huang2021gloria} & 2021 & CXR & ResNet-50 & --- & CheXpert reports & \cmark \\
CheXzero~\cite{tiu2022detection} & 2022 & CXR & ViT-B (CLIP) & --- & MIMIC-CXR pairs & \cmark \\
MedCLIP~\cite{wang2022medclip} & 2022 & CXR & Swin / ResNet & --- & unpaired images and texts & \cmark \\
BioViL~\cite{boecking2022most} & 2022 & CXR & ResNet-50 & --- & MIMIC-CXR pairs & \cmark \\
BiomedCLIP~\cite{zhang2023biomedclip} & 2023 & multi & ViT-B & --- & PMC-15M (15M pairs) & \cmark \\
CONCH~\cite{lu2024visuallanguage} & 2024 & pathology & ViT-B & --- & 1.17M image--caption pairs & \pmark \\
PLIP~\cite{huang2023visual} & 2023 & pathology & ViT-B (CLIP) & --- & OpenPath (Twitter) & \cmark \\
Mammo-CLIP~\cite{ghosh2024mammoclip} & 2024 & mammography & EfficientNet & --- & mammogram--report pairs & \cmark \\
CT-CLIP~\cite{hamamci2024ctrate} & 2024 & CT & 3-D ViT & --- & CT-RATE (25.7k volumes) & \cmark \\
\midrule
\multicolumn{7}{@{}l}{\emph{Multimodal LLMs (encoder--connector--LLM)}} \\
LLaVA-Med~\cite{li2023llavamed} & 2023 & multi & CLIP ViT-L & linear & PMC figure--caption + instructions & \cmark \\
Med-Flamingo~\cite{moor2023medflamingo} & 2023 & multi & CLIP ViT-L & resampler + x-attn & interleaved medical image--text & \cmark \\
MAIRA-2~\cite{bannur2024maira2} & 2024 & CXR & RAD-DINO & MLP & MIMIC-CXR, PadChest, USMix & \pmark \\
CheXagent~\cite{chen2024model} & 2024 & CXR & own CXR encoder & resampler & CheXinstruct & \cmark \\
LLaVA-Rad~\cite{chaves2024clinically} & 2024 & CXR & domain-adapted CLIP & MLP (adapter) & 697k image--text pairs & \cmark \\
RadFM~\cite{wu2023model} & 2023 & 2-D/3-D radiology & 3-D ViT & resampler & MedMD (16M scans) & \cmark \\
HuatuoGPT-Vision~\cite{chen2024huatuogptvision} & 2024 & multi & CLIP ViT-L & MLP & PubMedVision (1.3M VQA) & \cmark \\
Med-PaLM M~\cite{tu2023biomedical} & 2023 & multi & PaLM-E ViT & token fusion & MultiMedBench & \xmark \\
MedGemma~\cite{sellergren2025medgemma} & 2025 & multi & MedSigLIP & linear (pooled) & medical image--text + instructions & \pmark \\
Lingshu~\cite{team2025lingshu} & 2025 & multi & Qwen2.5-VL ViT & MLP & curated medical multimodal + text corpus & \cmark \\
\bottomrule
\end{tabular}
\end{table*}
